%==============================================================================
%  STAT 5200 - Applied Econometrics I
%  Homework 5
%
%  Compile with:  pdflatex "HW 5 - Report.tex"  (run twice for refs)
%
%  Directory layout assumed by this file:
%     HW 5 - Report.tex
%     R/        (R scripts)
%     figures/  (generated figures)
%     output/   (captured console output)
%==============================================================================
\documentclass[11pt]{article}

\usepackage[margin=1in]{geometry}
\usepackage[T1]{fontenc}
\usepackage[utf8]{inputenc}
\usepackage{amsmath,amssymb}
\usepackage{booktabs}
\usepackage{graphicx}
\usepackage{xcolor}
\usepackage{listings}
\usepackage[font=small,labelfont=bf]{caption}
\usepackage[hidelinks]{hyperref}

%------------------------------------------------------------------------------
%  Code presentation
%
%  Code is set as selectable text on a fixed character grid, source and console
%  output are given visually distinct styles, and source lines are kept to 76
%  characters so that nothing wraps.
%------------------------------------------------------------------------------
\definecolor{codekeyword}{RGB}{0,86,148}
\definecolor{codestring} {RGB}{163,21,21}
\definecolor{codecomment}{RGB}{106,115,125}
\definecolor{codebg}     {RGB}{248,248,250}
\definecolor{coderule}   {RGB}{220,220,226}

\lstdefinestyle{Rcode}{
  language        = R,
  basicstyle      = \ttfamily\footnotesize,
  keywordstyle    = \color{codekeyword}\bfseries,
  stringstyle     = \color{codestring},
  commentstyle    = \color{codecomment}\itshape,
  numbers         = left,
  numberstyle     = \ttfamily\scriptsize\color{codecomment},
  numbersep       = 9pt,
  backgroundcolor = \color{codebg},
  frame           = single,
  rulecolor       = \color{coderule},
  framesep        = 4pt,
  xleftmargin     = 15pt,
  xrightmargin    = 0pt,
  columns         = fixed,
  basewidth       = 0.5em,
  keepspaces      = true,
  tabsize         = 2,
  showstringspaces = false,
  breaklines      = true,
  breakatwhitespace = true,
  postbreak       = \mbox{\textcolor{codekeyword}{$\hookrightarrow$}\space},
  captionpos      = b,
  abovecaptionskip = 6pt,
}

\lstdefinestyle{Routput}{
  language        = {},
  basicstyle      = \ttfamily\footnotesize,
  keywordstyle    = \ttfamily,
  commentstyle    = \ttfamily,
  stringstyle     = \ttfamily,
  numbers         = none,
  backgroundcolor = \color{white},
  frame           = leftline,
  rulecolor       = \color{coderule},
  framesep        = 6pt,
  xleftmargin     = 20pt,
  breaklines      = true,
  captionpos      = b,
  abovecaptionskip = 6pt,
}

\lstset{style=Rcode}
\renewcommand{\lstlistingname}{Listing}

% Short inline code fragments.
\newcommand{\code}[1]{\texttt{#1}}

% Operators used throughout.
\DeclareMathOperator{\Var}{Var}
\DeclareMathOperator{\Cov}{Cov}
\DeclareMathOperator{\Avar}{Avar}
\DeclareMathOperator{\sd}{sd}
\DeclareMathOperator*{\plim}{plim}

% Placeholder marker. Every occurrence should be gone before submission;
% search the source for "\todo" to find what is still outstanding.
\newcommand{\todo}[1]{\textcolor{red}{\textbf{[TODO:} #1\textbf{]}}}

%------------------------------------------------------------------------------
\title{\vspace{-2em}Homework 5\\[0.3em]
       \large STAT 5200 --- Applied Econometrics I --- Fall 2026}
\author{Apostolos Kyratzis - 13130833}
\date{\today}

\begin{document}
\maketitle

% The handout says five questions in its introduction but numbers Q1--Q6.
% This template follows the numbered questions. No answers are completed yet.

\section{Q1: Reflection on Homework 3}
\label{sec:q1}
\subsection{Comparison with the solutions}
\todo{Compare your HW3 answers with the solutions and list any missed questions.}
\subsection{Areas for improvement}
\todo{Discuss what you could have done better.}

\section{Q2: Review}
\label{sec:q2}
% Read lab session notes 04 and 05 before working on the assignment.
% The handout requires no written answer for this question.
No written response required.

\section{Q3: Bootstrap and the delta method}
\label{sec:q3}
\subsection{(a) Nonparametric pairs bootstrap for HW3-Q3}
\todo{Describe the sampling, re-estimation, and inference procedure.}
\subsection{(b) Preferred method for future research}
\todo{State your preference and explain the reason.}

\section{Q4: Monte Carlo simulation and coverage rates}
\label{sec:q4}
% Use N = 100; B = 1000 for bootstrap methods. The handout permits a
% smaller Monte Carlo replication count, such as M = 500, if needed.
% Save scripts in R/, console results in output/, and plots in figures/.
\subsection{(a) Sample from the population}
\todo{Specify the data-generating process and generate the first sample.}
\subsection{(b) OLS estimation}
\todo{Compute the OLS coefficient estimates.}
\subsection{(c) Asymptotic confidence interval}
\todo{Compute the 95\% interval for the slope and its coverage indicator.}
\subsection{(d) Empirical coverage under homoscedasticity}
\todo{Repeat the experiment and discuss the empirical coverage rate.}
\subsection{(e) Empirical coverage under heteroscedasticity}
\todo{Repeat with the modified error term and compare coverage rates.}
\subsection{(f) Heteroscedasticity-consistent confidence intervals}
\todo{Report and discuss coverage for HC0, HC1, HC2, and HC3.}
\subsection{(g) Bootstrap and jackknife confidence intervals}
\todo{Report and discuss coverage for pairs, residual, and wild bootstrap,
and for the jackknife.}

\section{Q5: Proxy variables and the return to education}
\label{sec:q5}
\subsection{(a) Estimates with and without IQ}
\todo{Compare the parameter estimates reported in lab05.}
\subsection{(b) Requirements for a valid proxy}
\todo{State the two requirements and discuss whether IQ satisfies them.}
\subsection{(c) Sign of the education coefficient in the ability model}
\todo{Infer the sign of the coefficient from the empirical results and explain.}

\section{Q6: Bootstrapping in an empirical paper}
\label{sec:q6}
\subsection{(1--2) Paper selection and citation}
\todo{Find a qualifying published empirical paper and give its full citation.}
\subsection{(3) Research question}
\todo{Summarize the main research question in one paragraph.}
\subsection{(4) Empirical model}
\todo{Describe the model, outcome, regressors, number of variables, and sample size.}
\subsection{(5) Resampling method}
\todo{Identify the bootstrap or jackknife method used.}
\subsection{(6) Motivation for the method}
\todo{Explain the authors' motivation and assess whether it is convincing.}
\subsection{(7) Robustness to heteroscedasticity}
\todo{Discuss whether the method is robust to heteroscedasticity.}

% Add code and output listings with \lstinputlisting when scripts exist.
% Example (uncomment after creating the file):
% \appendix
% \section{R code}
% \lstinputlisting[style=Rcode]{R/q4_coverage_rate.R}

\end{document}
